% Generated draft: requires mathematical and source review.
\documentclass[11pt]{article}
\usepackage[margin=25mm]{geometry}
\usepackage{fontspec}
\setmainfont{Latin Modern Roman}
\usepackage{amsmath,amssymb,mathtools,unicode-math}
\setmathfont{Latin Modern Math}
\usepackage{enumitem,needspace,xurl,hyperref}
\setlength{\parindent}{0pt}
\setlength{\parskip}{5pt}
\setlength{\emergencystretch}{3em}
\newcommand{\N}{\mathbb{N}}
\newcommand{\Z}{\mathbb{Z}}
\newcommand{\Q}{\mathbb{Q}}
\newcommand{\R}{\mathbb{R}}
\newcommand{\C}{\mathbb{C}}
\newcommand{\F}{\mathbb{F}}
\newcommand{\A}{\mathbb{A}}
\newcommand{\PP}{\mathbb{P}}
\newcommand{\E}{\mathbb{E}}
\newcommand{\eps}{\varepsilon}
\newcommand{\dd}{\,\mathrm{d}}
\begin{document}
\section*{UnsolvedMath ID: 20003172}
\textbf{Source problem code:} catalogue entry
\subsection{Definitions and mathematical statement}
Use the usual meaning of quasi-positive sectional curvature: $\sec\geq0$ everywhere and all two-planes at some point have strictly positive curvature. For a Riemannian manifold $(M,g)$ put $\operatorname{cohom}(M,g)=\dim(M/\operatorname{Iso}(M,g))$, using the full isometry group and its quotient metric. The source asks, for each positive integer $k$, whether there is an integer $n_0(k)$ such that every intended quasi-positively curved $(M^n,g)$ with $n\geq n_0(k)$ and cohomogeneity at most $k$ admits a tower
\[
M_0=M^n\subset M_1^{n+k}\subset M_2^{n+2k}\subset\cdots
\]
with isometric totally geodesic inclusions, quasi-positive metrics $g_i$, a direct union modeling $BG$ for some Lie group $G$, and
\[
M_i/\operatorname{Iso}(M_i,g_i)\cong M_0/\operatorname{Iso}(M_0,g_0)
\]
isometrically for every $i$. Compactness, completeness, the topology on the union, and whether modeling $BG$ means homotopy equivalence are not explicit in the source. The model statements below use closed connected manifolds and the standard weak/direct-limit topology; a metric completion is a different construction. The codimension is exactly $k$ at every step, not merely bounded by $k$.

The aid reports full model towers with standard metrics: round spheres $S^{n+ik}$ and real projective spaces $\R\mathrm P^{n+ik}$ for any positive $k$; complex projective spaces $\C\mathrm P^{a+i(k/2)}$ when $k$ is even and $n=2a$; and quaternionic projective spaces $\mathbb H\mathrm P^{a+i(k/4)}$ when $4\mid k$ and $n=4a$. Their full isometry groups are transitive, so all quotient spaces are points. Their stable unions have respectively the homotopy types of $B\{1\}$, $B(\Z/2\Z)$, $BU(1)$, and $BSp(1)$. The sphere union is contractible. These individual families do not establish the universal threshold $n_0(k)$.

The aid also invokes Wilking's real, complex, and quaternionic projectivized tangent-bundle families as near-solutions, with fixed steps $k=2,4,8$ respectively. Its stable flag description uses the unit scalar groups $O(1)=\Z/2\Z$, $U(1)$, and $Sp(1)$, so the corresponding limits are
\[
B((\Z/2\Z)^2),\qquad B(U(1)^2),\qquad B(Sp(1)^2).
\]
For those chosen families and metrics it reports fixed dimension steps, quasi-positive curvature, totally geodesic inclusions, the cohomogeneity bounds, and the stable classifying-space clauses. The finite-stage metrics and embedding formulas are not supplied here. Crucially, stage-independence of the quotient metric by the \emph{full} isometry group remains uncertified; constancy for a chosen subgroup action would not prove that clause.

Finally the aid reports a conditional connectedness argument requiring a normal-geodesic visibility hypothesis to transfer Wilking's positive-curvature stabilization method. The precise hypothesis and proposition are absent, so it supplies no general stabilization theorem for quasi-positive curvature. The full quantified source question remains unresolved by the reported model and near-model towers.
\subsection{Short English statement}
The source asks for a fixed-codimension stabilization tower for every sufficiently large low-cohomogeneity quasi-positively curved manifold, preserving the full-isometry quotient metric and ending in a classifying space. Standard sphere and projective towers satisfy all clauses in their own families. The supplied Wilking near-models retain a specific gap: their full-isometry quotient metrics are not certified constant. Neither family proves the universal threshold statement.
\subsection{Sources}
\begin{itemize}
\item Local registry: permanent catalogue entry, exact source code catalogue entry; \texttt{aim-topology-notes.json}, \texttt{notes[259]}.
\item AIM workshop \emph{Manifolds with nonnegative sectional curvature}, Problem 47, \url{https://aimath.org/WWN/nnsectcurvature/nnsectcurvature.pdf}. No proposer, proposal year, remarks, or original literature reference is supplied.
\item Verification reports checking page 4 and its PDF content-stream placement in the inherited run, repairing the display to $M_0=M^n\subset M_1^{n+k}\subset M_2^{n+2k}\subset\cdots$. Stray greater-than signs and the following section heading ``7. Ricci Flow'' are extraction artifacts, not mathematical conditions.
\item Supplied research attempt 1, registry updated 14 August 2026: partial result, open-at-run label, low novelty confidence, machine-generated and not peer reviewed. It names Wilking's projectivized tangent-bundle and positive-curvature stabilization work without full titles, theorem numbers, or authoritative citations, and explicitly leaves the full-isometry quotient clause uncertified.
\end{itemize}
\end{document}
